\section{Discussion}
\label{sec:discussion}

The paper is easiest to read if its two outputs are kept separate. Theorems
calibrate fixed-node Stage~A screening under a complete permutation null. The
experiments evaluate the full tree or forest as an adaptive ranking procedure.
The same algorithm can realize both regimes, but only the first carries the
finite-sample guarantees stated in \cref{sec:theory}.

The structural lesson from the empirical sections is a strictness--depth
tradeoff. Conservative nodewise screening protects against spurious rejections,
yet in forests it also suppresses depth, which is the main mechanism by which
weaker or redundant features enter the ranking. That is why the method can
look strong on structured problems but lose ground in weak-signal or very
high-dimensional regimes: the cost of inferential discipline is paid through
shallower trees.

Exchangeability remains the key modeling assumption. If labels are clustered,
longitudinal, blocked, or otherwise dependent, naive permutations do not
preserve the null distribution. The same caution applies to root resampling:
any bootstrap design that breaks the relevant exchangeability should be read as
an engineering choice rather than as part of the calibration story.

Ranking stability is another practical consideration. As with other
forest-based importance methods, the exact top-$k$ set can vary across seeds
even when the broader ranking is fairly stable. Applications that require a
small, reproducible feature set may therefore prefer a simpler screening rule
or an explicit stability-selection wrapper.

The present manuscript is organized to keep the paper robust to future result
updates: the main claims attach to the fixed-node theory, the benchmark design,
and the strictness--depth tradeoff rather than to any single ranking table.
That is also the natural path forward methodologically. Less conservative
nodewise corrections, sample-split treatments of Stage~B, and stronger
stability mechanisms are all promising ways to widen the range of problems on
which conditional-inference forests are attractive.
